%==========================================================
%  CHAPITRE 10 — ANNEXES
%==========================================================
\chapter{Annexes}
\label{chap:annexes}

%──────────────────────────────────────────────────────────
\section{Annexe A — Cahier des charges détaillé}
\label{ann:cdc}
%──────────────────────────────────────────────────────────

\todo{Insérer ici le cahier des charges complet fourni par le maître
de stage, ou y faire référence si trop volumineux.}

\placeholder{Cahier des charges complet (document PDF joint ou extrait
— à insérer ici si applicable)}

%──────────────────────────────────────────────────────────
\section{Annexe B — Glossaire}
\label{ann:glossaire}
%──────────────────────────────────────────────────────────

\begin{description}[leftmargin=3cm, style=sameline]
  \item[\textbf{API}]  \textit{Application Programming Interface} —
        Interface permettant à deux logiciels de communiquer.
  \item[\textbf{CI/CD}] \textit{Continuous Integration / Continuous Delivery} —
        Pratiques DevOps automatisant les phases de build, test et déploiement.
  \item[\textbf{CRUD}] \textit{Create, Read, Update, Delete} —
        Les quatre opérations de base sur une base de données.
  \item[\textbf{DTO}]  \textit{Data Transfer Object} —
        Objet utilisé pour transporter des données entre les couches
        d'une application.
  \item[\textbf{JWT}]  \textit{JSON Web Token} —
        Standard ouvert pour l'authentification sans état.
  \item[\textbf{MVC}]  \textit{Model-View-Controller} —
        Patron architectural séparant la logique métier, la vue et le contrôleur.
  \item[\textbf{ORM}]  \textit{Object-Relational Mapping} —
        Technique de correspondance entre objets et tables relationnelles.
  \item[\textbf{REST}] \textit{Representational State Transfer} —
        Style architectural pour les services web.
  \item[\textbf{Scrum}] Cadre Agile de développement itératif basé sur des sprints.
  \item[\textbf{SLA}]  \textit{Service Level Agreement} —
        Accord définissant le niveau de service attendu.
  \item[\textbf{UML}]  \textit{Unified Modeling Language} —
        Langage de modélisation graphique standardisé.
  \item[\todo{Terme}]  \todo{Définition}
\end{description}

%──────────────────────────────────────────────────────────
\section{Annexe C — Documentation de l'API}
\label{ann:api}
%──────────────────────────────────────────────────────────

\todo{Insérer ici la documentation Swagger / OpenAPI des endpoints
développés, ou un tableau récapitulatif des routes.}

\begin{center}
\begin{tabular}{>{\ttfamily\footnotesize}l >{\sffamily\footnotesize\bfseries}c
                >{\footnotesize}l >{\footnotesize}l}
  \toprule
  \textbf{Route}              & \textbf{Méthode} & \textbf{Description}
                                                  & \textbf{Auth} \\
  \midrule
  /api/v1/auth/login          & POST   & Authentification   & Non \\
  /api/v1/auth/refresh        & POST   & Rafraîchir token   & Oui \\
  /api/v1/\todo{resource}     & GET    & Lister resources   & Oui \\
  /api/v1/\todo{resource}     & POST   & Créer resource     & Oui \\
  /api/v1/\todo{resource}/\{id\}& GET  & Lire resource      & Oui \\
  /api/v1/\todo{resource}/\{id\}& PUT  & Modifier resource  & Oui \\
  /api/v1/\todo{resource}/\{id\}& DELETE& Supprimer resource & Admin \\
  \todo{Autres routes}        & \todo{}& \todo{}           & \todo{} \\
  \bottomrule
\end{tabular}
\end{center}

%──────────────────────────────────────────────────────────
\section{Annexe D — Guide d'installation}
\label{ann:install}
%──────────────────────────────────────────────────────────

\subsection{Prérequis}

\begin{itemize}
  \item \todo{Java 17+ / Python 3.11+ / Node.js 18+}
  \item \todo{Docker 24.0+} et Docker Compose v2
  \item \todo{Git 2.40+}
  \item \todo{Autre prérequis}
\end{itemize}

\subsection{Étapes d'installation}

\begin{lstlisting}[
  language={},
  caption={Commandes d'installation et de lancement},
  label={lst:install}
]
# 1. Cloner le depot
git clone https://github.com/[A COMPLETER]/[projet].git
cd [projet]

# 2. Copier la configuration
cp .env.example .env
# Editer .env avec vos valeurs (DB, secrets, etc.)

# 3. Lancer avec Docker Compose
docker compose up --build -d

# 4. Initialiser la base de donnees (si applicable)
docker compose exec backend [commande de migration]

# 5. Acceder a l'application
# Frontend : http://localhost:3000
# API :      http://localhost:8080/api/v1
# Swagger :  http://localhost:8080/swagger-ui
\end{lstlisting}

%──────────────────────────────────────────────────────────
\section{Annexe E — Planning Gantt}
\label{ann:gantt}
%──────────────────────────────────────────────────────────

\placeholder{Diagramme de Gantt du projet — figures/gantt.png\\
(à générer avec MS Project, Excel, GanttProject ou pgfgantt)}

\todo{Commenter le Gantt : jalons importants, tâches critiques,
marges, décalages par rapport au plan initial.}
